\begin{helper}
Let \(\ell>0\).  If there is an \(\ell\)-coloring
\[
  \chi:E(K_n)\to[\ell]
\]
with no monochromatic clique of size \(s\), in the sense of
\(\noderef{NoMonochromaticCliqueColoring}\), then
\[
  n<\noderef{MulticolorRamseyNumber}(s,\ell).
\]
\end{helper}
\begin{proof}
Let \(\chi\) be such a coloring.  The definition
\(\noderef{MulticolorRamseyProperty}(s,\ell,n)\) says that every
\(\ell\)-coloring of \(K_n\) has a monochromatic \(s\)-clique.  Applying this
property to the particular coloring \(\chi\) would contradict
\(\noderef{NoMonochromaticCliqueColoring}\).  Hence
\(\noderef{MulticolorRamseyProperty}(s,\ell,n)\) is false.

By \(\noderef{MulticolorRamseyPropertyNonempty}\), the set
\[
  A=\{m:\noderef{MulticolorRamseyProperty}(s,\ell,m)\}
\]
is nonempty.  Therefore the natural-number infimum defining
\(\noderef{MulticolorRamseyNumber}(s,\ell)\) is a member of \(A\).  Suppose,
for contradiction, that
\[
  \noderef{MulticolorRamseyNumber}(s,\ell)\le n.
\]
Since the Ramsey property is monotone by
\(\noderef{MulticolorRamseyPropertyMonotone}\), membership of the infimum in
\(A\) implies \(\noderef{MulticolorRamseyProperty}(s,\ell,n)\).  This
contradicts the previous paragraph.  Hence the assumed inequality is false,
and so \(n<\noderef{MulticolorRamseyNumber}(s,\ell)\).
\end{proof}

